%% =====================================================================
%%  Статья для журнала «Искусственный интеллект и принятие решений» (ИИиПР)
%%  Paper 1 — GigaEmbeddings
%%  Формат ИИиПР: RU+EN единым файлом, <=20 c. A4, Times New Roman 11pt,
%%  интервал 1.5, поля 3.5/3.5/2.5/2.5, аннотация 100-120 слов,
%%  ключевые слова <=8-10, литература <=50 (ГОСТ Р 7.0.5-2008),
%%  русскоязычные источники перед иноязычными. Итог сдаётся в MS Word.
%% =====================================================================
\documentclass[12pt,a4paper]{article}

\usepackage[T2A]{fontenc}
\usepackage[utf8]{inputenc}
\usepackage[english,russian]{babel}
\usepackage{geometry}
\geometry{a4paper,top=3.5cm,bottom=3.5cm,left=2.5cm,right=2.5cm}
\usepackage{amsmath,amssymb,amsfonts,mathtools}
\usepackage{graphicx}
\usepackage{booktabs}
\usepackage{array}
\usepackage{setspace}
\usepackage{tikz}
\usetikzlibrary{positioning,shapes.geometric,arrows.meta,fit,calc}
\onehalfspacing

\begin{document}

%% ---------------------------------------------------------------
%%  РУССКАЯ ЧАСТЬ
%% ---------------------------------------------------------------
\noindent\textbf{УДК 004.852}

\bigskip
\begin{center}
  {\large\bfseries Иерархическое инструктивное дообучение декодерной\\
  языковой модели для русскоязычных текстовых эмбеддингов}

  \medskip
  {\bfseries Колодин Е.\,И.$^{1}$, Янина А.\,О.$^{1}$} % [ДОБАВИТЬ соавторов команды GigaEmbeddings при необходимости]

  \smallskip
  {\itshape $^{1}$ Московский физико-технический институт\\
  (национальный исследовательский университет), г. Долгопрудный}

  \smallskip
  {\small ekolodin98@gmail.com}
\end{center}

\medskip
\noindent\textbf{Аннотация.} % 100-120 слов, один абзац, без формул и сокращений
Представлен фреймворк GigaEmbeddings для построения высокопроизводительных
русскоязычных текстовых эмбеддингов на основе декодерной языковой модели.
Предложено иерархическое трёхэтапное инструктивное дообучение: адаптация
декодера к двунаправленному вниманию и латентному пулингу, контрастивное
предобучение на корпусах веб-масштаба, дообучение на жёстких негативах и
мультизадачная генерализация. Рассмотрена интеграция модели в двухэтапный
конвейер информационного поиска с приближённым поиском ближайших соседей и
проанализирована вычислительная эффективность прунинга глубоких слоёв при
промышленном развёртывании. На бенчмарке MTEB(rus, v1.1) базовая версия модели
достигает первого места со средним баллом 74,16, превосходя существенно более
крупные многоязычные и русскоязычные базовые модели. Дальнейшее развитие модели~---
упрощение архитектуры для совместимости с движками быстрого инференса и обучение
со слиянием весов~--- повышает средний балл обновлённой версии до 74,57 на
русскоязычной дорожке и распространяет модель на английскую, кодовую и
многоязычную дорожки MTEB. Абляционные исследования подтверждают вклад каждого
этапа обучения.

\medskip
\noindent\textbf{Ключевые слова:} текстовые эмбеддинги, информационный поиск,
большие языковые модели, контрастивное обучение, инструктивное дообучение,
приближённый поиск ближайших соседей, русский язык, MTEB.

%% ===============================================================
\section{Введение}
%% ===============================================================
Текстовые эмбеддинги~--- векторные представления, кодирующие семантику
естественного языка,~--- лежат в основе информационного поиска,
вопросно-ответных систем, оценки семантического сходства и систем генерации с
извлечением информации (retrieval-augmented generation, RAG). Качество
эмбеддингов напрямую определяет полноту и точность первичного этапа извлечения,
который в промышленных системах реализуется приближённым поиском ближайших
соседей по индексу из миллионов документов.

Для русского языка построение универсальной модели эмбеддингов осложнено
ограниченностью размеченных данных и морфологическим богатством языка.
В настоящей работе показано, что декодерная языковая модель, прошедшая
иерархическое инструктивное дообучение, обеспечивает качество эмбеддингов уровня
state-of-the-art при умеренном числе параметров. Базовая методика
предварительно представлена в~[\ref{ref:acl}]; настоящая работа существенно
расширяет её по трём направлениям. Во-первых, формализована интеграция модели в
двухэтапный поисковый конвейер с приближённым поиском ближайших соседей и
аналитически оценена вычислительная сложность приближённого поиска
(раздел~\ref{sec:ann}). Во-вторых, проанализирована эффективность промышленного
развёртывания усечённой модели (раздел~\ref{sec:deploy}). В-третьих, описано
дальнейшее развитие модели~--- упрощение архитектуры и обучение со слиянием
весов,~--- повышающее качество и распространяющее модель на английскую, кодовую
и многоязычную дорожки бенчмарка MTEB (раздел~\ref{sec:v2}).

\textbf{Вклад работы.} (1)~Предложен трёхэтапный конвейер инструктивного
дообучения декодерной модели для эмбеддингов. (2)~Формализована интеграция
модели в двухэтапный поиск и аналитически оценена сложность приближённого
поиска. (3)~Получено первое место на MTEB(rus, v1.1) со средним баллом 74,16.
(4)~Описаны упрощение архитектуры и обучение со слиянием весов, повышающие
качество до 74,57 и распространяющие модель на английскую, кодовую и
многоязычную дорожки. (5)~Абляционные исследования количественно разделяют
вклад этапов обучения, прунинга и стратегии промптирования.

%% ===============================================================
\section{Архитектура модели}\label{sec:arch}
%% ===============================================================
\subsection{Базовая модель и двунаправленное внимание}
В качестве основы используется предобученный чекпоинт декодерной языковой модели
GigaChat-3B~[\ref{ref:gigachat}]. Масштаб 3~млрд параметров отражает баланс
между вычислительной эффективностью и репрезентативной мощностью. На этапе
обучения каузальные маски внимания удаляются, что обеспечивает двунаправленное
контекстное моделирование: каждый токен учитывает как предшествующие, так и
последующие позиции~--- условие, критическое для качества векторного
представления последовательности.

\subsection{Прунинг глубоких слоёв}
Следуя наблюдению о низкой значимости глубоких слоёв трансформера~[\ref{ref:gromov}],
исследован прунинг 25\,\% блоков: из 36 блоков (модули самовнимания и
feed-forward) удаляются 9 наиболее глубоких по индексу, что даёт компактную
модель приблизительно на 2{,}5~млрд параметров. Удаляются только блоки
трансформера, тогда как таблица эмбеддингов токенов и голова пулинга
сохраняются; этим объясняется, что 25\,\% удалённых блоков соответствуют
сокращению общего числа параметров примерно на 17\,\%. Полная модель выпущена в
открытый доступ, усечённая применяется в промышленном развёртывании для снижения
задержки инференса (раздел~\ref{sec:deploy}).

\subsection{Латентный пулинг внимания}
Агрегация скрытых активаций декодера в финальный вектор выполняется головой
латентного пулинга внимания (latent attention pooling) в духе
подхода~[\ref{ref:nvembed}]. Механизм реализован как кросс-внимание: выход
декодера выступает запросом $\mathbf{Q}$, а обучаемый латентный массив~---
ключом $\mathbf{K}$ и значением $\mathbf{V}$; результат обрабатывается
многослойным перцептроном.

%% ===============================================================
\section{Конвейер обучения}\label{sec:training}
%% ===============================================================
Предложен иерархический трёхэтапный конвейер инструктивного дообучения.
Во всех источниках выполняется дедупликация, а разбиение на обучающую и
оценочную части проводится на уровне документов для предотвращения утечки
данных бенчмарка.

\subsection{Функция потерь}
Используется функция потерь InfoNCE~[\ref{ref:infonce}] с фиксированной
температурой $\tau = 0{,}02$:
\begin{equation}\label{eq:infonce}
  \mathcal{L} = -\log
  \frac{\varphi(q, d^{+})}{\varphi(q, d^{+}) + \sum_{n_i \in \mathcal{N}} \varphi(q, n_i)},
\end{equation}
где $\mathcal{N}$~--- множество внутрибатчевых, кросс-батчевых и жёстких
негативов, а функция совпадения запроса $q$ и позитивного документа $d^{+}$
определяется через косинусную близость эмбеддингов $\mathbf{h}_q,\mathbf{h}_{d^{+}}$:
\begin{equation}\label{eq:match}
  \varphi(q, d^{+}) = \exp\!\left(\tfrac{1}{\tau}\,\cos(\mathbf{h}_q, \mathbf{h}_{d^{+}})\right).
\end{equation}

\subsection{Этап 1: контрастивное предобучение}
Используются корпуса веб-масштаба в формате «заголовок--пассаж» (Wikipedia,
Reddit, StackExchange и научный корпус S2ORC~[\ref{ref:s2orc}]) и добытые
пассажи; для разнообразия запросов применяется синтетическая генерация
инструктивной моделью. Обучение ведётся с внутрибатчевыми и кросс-батчевыми
негативами при большом размере батча (16\,384), что максимизирует разнообразие
негативов при адаптации декодера к двунаправленному вниманию.

\subsection{Этап 2: дообучение на жёстких негативах}
На высококачественных размеченных IR-датасетах (MS~MARCO~[\ref{ref:msmarco}],
Natural Questions~[\ref{ref:nq}], SQuAD~[\ref{ref:squad}],
MIRACL~[\ref{ref:miracl}], Mr.~TyDi~[\ref{ref:mrtydi}]) функция~(\ref{eq:infonce})
расширяется семью курируемыми жёсткими негативами на запрос. Жёсткие негативы
добываются предобученной моделью как кандидаты с высоким сходством, из которых
порогом семантического сходства отсекаются вероятные ложноотрицательные
примеры.

\subsection{Этап 3: мультизадачная генерализация}
В обучающую смесь добавляются задачи классификации и кластеризации. Для
предотвращения ложных негативов в нерелевантных задачах внутрибатчевые негативы
удаляются; применяется задачно-специфичное инструктивное дообучение с единой
функцией потерь.

\subsection{Гиперпараметры}
Гиперпараметры этапов приведены в табл.~\ref{tab:hp}.
\begin{table}[ht]
  \centering
  \caption{Гиперпараметры этапов обучения}\label{tab:hp}
  \begin{tabular}{lccc}
    \toprule
    Параметр & Этап 1 & Этап 2 & Этап 3 \\
    \midrule
    Размер батча                & 16\,384 & 512 & 512 \\
    Макс. длина, токенов        & 512 & 512 & 512 \\
    Темп обучения               & $1{\cdot}10^{-5}$ & $1{\cdot}10^{-5}$ & $1{\cdot}10^{-5}$ \\
    Шаги разогрева              & 1000 & 500 & 500 \\
    Длительность обучения       & 6000 шагов & 3 эпохи & 3 эпохи \\
    Затухание весов             & 0{,}01 & 0{,}01 & 0{,}01 \\
    Жёсткие негативы на запрос  & ---  & 7 & 7 \\
    Внутрибатчевые негативы     & да   & да & нет \\
    Температура $\tau$          & 0{,}02 & 0{,}02 & 0{,}02 \\
    Оптимизатор                 & \multicolumn{3}{c}{AdamW} \\
    \bottomrule
  \end{tabular}
\end{table}

\subsection{Стратегия промптирования}
Сравнены два способа: \textit{prefix} (фиксированный префикс) и \textit{instruct}
(полная инструкция о задаче). Инструктивное промптирование превосходит
префиксное на 0{,}8 балла ruMTEB (раздел~\ref{sec:exp}).

\subsection{Воспроизводимость}
Финальный вектор последовательности формируется головой латентного пулинга
внимания по скрытым активациям последнего слоя (раздел~\ref{sec:arch});
размерность итогового эмбеддинга составляет 2048, максимальная длина входа~---
512 токенов. Для предотвращения утечки бенчмарка дедупликация и разбиение
проводятся на уровне документов. Полная и усечённая версии модели, а также
конфигурации обучения выпущены в открытый доступ вместе с базовой
публикацией~[\ref{ref:acl}], что обеспечивает точное воспроизведение архитектуры
и гиперпараметров.

%% ===============================================================
\section{Интеграция в двухэтапный поиск и приближённый поиск
ближайших соседей}\label{sec:ann}
%% ===============================================================
В промышленном контуре эмбеддинги обслуживают первичный этап извлечения, к
которому предъявляются жёсткие требования по задержке. Структура двухэтапного
поискового конвейера показана на рис.~\ref{fig:two-stage}: bi-энкодер
GigaEmbeddings отбирает кандидатов по индексу документов, после чего
кросс-энкодер уточняет ранжирование. Точный поиск ближайших соседей по $N$
документам имеет сложность $O(Nd)$ на запрос ($d = 2048$~--- размерность эмбеддинга) и
непрактичен при больших $N$. Поэтому применяется приближённый поиск ближайших
соседей (ANN), реализуемый двумя взаимодополняющими семействами индексов.

\begin{figure}[ht]
  \centering
  \resizebox{\textwidth}{!}{%
    \begin{tikzpicture}[
    font=\sffamily,
    node distance=1.3cm and 1.9cm,
    box/.style={
        draw,
        rounded corners,
        minimum width=2.8cm,
        minimum height=1.1cm,
        align=center,
        thick
    },
    db/.style={
        draw,
        cylinder,
        shape border rotate=90,
        aspect=0.3,
        minimum width=2.2cm,
        minimum height=1.6cm,
        align=center,
        thick
    },
    stage/.style={
        draw,
        dashed,
        rounded corners,
        inner sep=0.35cm
    },
    arr/.style={-{Stealth[length=3mm]}, thick}
]
    %% Запрос
    \node[box] (query) {Запрос\\пользователя};
    %% Стадия 1: отбор кандидатов
    \node[box, right=of query] (retriever) {Bi-encoder\\(GigaEmbeddings)};
    \node[db, above=of retriever] (index) {Индекс\\документов};
    %% Стадия 2: переранжирование
    \node[box, right=of retriever] (reranker) {Cross-encoder\\реранкер};
    %% Результат
    \node[box, right=of reranker] (result) {Top-$n$\\документов};

    %% Подписи стадий
    \node[stage, fit=(retriever)(index), label={[font=\sffamily\small]above:Отбор (recall)}] {};
    \node[stage, fit=(reranker), label={[font=\sffamily\small]above:Уточнение (precision)}] {};

    %% Стрелки
    \draw[arr] (query) -- (retriever);
    \draw[arr] (index) -- (retriever);
    \draw[arr] (retriever) -- node[above, font=\small] {top-$k$} (reranker);
    \draw[arr] (reranker) -- (result);
\end{tikzpicture}
}
  \caption{Двухэтапный поисковый конвейер: bi-энкодер (GigaEmbeddings) отбирает
  top-$k$ кандидатов по индексу, кросс-энкодер уточняет top-$n$.}\label{fig:two-stage}
\end{figure}

\textbf{Графовый индекс (HNSW).} Иерархический навигируемый малый
мир~[\ref{ref:hnsw}] строит многоуровневый граф близости; жадный спуск по
уровням обеспечивает ожидаемую сложность поиска $O(\log N)$ при контролируемой
полноте, регулируемой параметрами степени вершины и ширины списка кандидатов.

\textbf{Инвертированный индекс с квантованием (IVF--PQ)~[\ref{ref:faiss}].} Пространство разбивается
на $n_{\text{list}}$ ячеек Вороного грубым квантователем; на запросе
просматриваются лишь $n_{\text{probe}}$ ближайших ячеек. Произведённое
квантование (product quantization, PQ) делит вектор на $M$ подвекторов, каждый
из которых кодируется ближайшим центроидом из словаря размера $k^{*}$, что даёт
сжатие до $M\log_2 k^{*}$ бит на вектор. Расстояние аппроксимируется суммой
предвычисленных расстояний по подпространствам:
\begin{equation}\label{eq:pq}
  \widehat{\lVert \mathbf{q}-\mathbf{x}\rVert}^{2}
  = \sum_{m=1}^{M} \big\lVert \mathbf{q}^{(m)} - \mathbf{c}^{(m)}_{j_m(\mathbf{x})} \big\rVert^{2},
\end{equation}
где $\mathbf{c}^{(m)}_{j}$~--- центроиды $m$-го подпространства. Полнота поиска
управляется параметрами $n_{\text{probe}}$, $M$ и $k^{*}$, задающими компромисс
между задержкой, памятью и качеством.

\textbf{Двухэтапный поиск.} ANN-индекс формирует множество кандидатов, которое
уточняется кросс-энкодером (переранжированием). Эмбеддинги GigaEmbeddings служат
первым этапом; качество переранжирования (категория Rerank в табл.~\ref{tab:res})
характеризует пригодность модели в этой роли.

\textbf{Компромисс полноты и задержки.} Оба семейства индексов задают
монотонный компромисс между полнотой $\mathrm{Recall}@k$ и задержкой,
управляемый своими параметрами. Для HNSW увеличение ширины списка кандидатов
$\mathrm{efSearch}$ и степени вершины повышает полноту ценой роста числа
просматриваемых вершин и, следовательно, задержки $O(\log N)$ с большей
константой. Для IVF--PQ полнота растёт с числом просматриваемых ячеек
$n_{\text{probe}}$ (вплоть до $O(N)$ при $n_{\text{probe}} = n_{\text{list}}$) и
с точностью квантования $M, k^{*}$, что одновременно увеличивает память
$M\log_2 k^{*}$ бит на вектор и время оценки расстояния~(\ref{eq:pq}). Графовые
индексы обычно обеспечивают более высокую полноту при малой задержке, тогда как
IVF--PQ предпочтителен при жёстких ограничениях по памяти на корпусах в сотни
миллионов векторов. Поскольку эмбеддинг фиксированной размерности $d$
вычисляется моделью однократно на запрос, выбор индекса не затрагивает саму
модель и определяется эксплуатационными ограничениями развёртывания;
количественная калибровка кривой полнота--задержка на конкретном корпусе и
оборудовании составляет предмет инженерной настройки.

%% ===============================================================
\section{Экспериментальная оценка}\label{sec:exp}
%% ===============================================================
\subsection{Постановка}
Качество оценивается на русскоязычном бенчмарке MTEB(rus, v1.1)~[\ref{ref:mteb}],
охватывающем семь категорий задач. Модель сравнивается с сильными базовыми
моделями: Qwen3-Embedding-8B~[\ref{ref:qwen}], FRIDA, multilingual-E5-large-instruct,
E5-Mistral-7B-instruct~[\ref{ref:e5mistral}], GritLM-7B~[\ref{ref:gritlm}],
BGE-M3~[\ref{ref:bgem3}].

\subsection{Основные результаты}
По данным публичного рейтинга модель Giga-Embeddings-instruct занимает первое
место со средним баллом 74,16, опережая Qwen3-Embedding-8B (71,42) и FRIDA
(70,95); при этом модель более чем вдвое компактнее ближайшего конкурента
Qwen3-Embedding-8B (3,23 против 8,0~млрд параметров). Результаты по категориям и
число параметров приведены в табл.~\ref{tab:res}.

\begin{table}[ht]
  \centering
  \caption{Сравнение на MTEB(rus, v1.1) по семи категориям задач. Число
  параметров приведено приближённо по карточкам моделей. Обозначения:
  Cls~--- классификация, Clust~--- кластеризация, MLCls~--- мультиклассовая
  классификация, PairCls~--- классификация пар, Rerank~--- переранжирование,
  Retr~--- информационный поиск}\label{tab:res}
  \resizebox{\textwidth}{!}{%
  \begin{tabular}{lccccccccc}
    \toprule
    Модель & Парам., млрд & Среднее & Cls & Clust & MLCls & PairCls & Rerank & Retr & STS \\
    \midrule
    Giga-Embeddings-instruct       & 3,23 & \textbf{74,16} & \textbf{76,66} & 67,48 & \textbf{57,07} & \textbf{79,57} & \textbf{74,02} & \textbf{81,41} & 75,75 \\
    Qwen3-Embedding-8B             & 8,0  & 71,42 & 74,35 & 66,72 & 44,58 & 71,33 & 70,19 & 78,19 & \textbf{79,28} \\
    FRIDA                          & 0,82 & 70,95 & 73,71 & \textbf{67,49} & 52,23 & 66,45 & 71,50 & 76,38 & 74,35 \\
    E5-Mistral-7B-instruct         & 7,1  & 67,04 & 69,07 & 64,24 & 42,93 & 60,81 & 69,96 & 73,29 & 73,71 \\
    GritLM-7B                      & 7,2  & 65,02 & 69,92 & 64,30 & 41,96 & 58,93 & 69,99 & 55,52 & 74,63 \\
    multilingual-E5-large-instruct & 0,56 & 62,41 & 66,28 & 63,13 & 41,15 & 63,89 & 64,35 & 48,45 & 76,48 \\
    BGE-M3                         & 0,57 & 61,63 & 60,44 & 52,38 & 34,86 & 60,60 & 69,71 & 75,20 & 73,68 \\
    \bottomrule
  \end{tabular}}
\end{table}

\subsection{Абляционные исследования}
Таблица~\ref{tab:abl} разделяет вклад компонентов. Контрастивное предобучение
даёт +0{,}6 балла (прежде всего за счёт информационного поиска), а инструктивное
промптирование~--- +0{,}8 балла относительно префиксного.

\begin{table}[ht]
  \centering
  \caption{Абляционные исследования (средний балл ruMTEB по 23 задачам, по
  данным базовой публикации~[\ref{ref:acl}])}\label{tab:abl}
  \begin{tabular}{lc}
    \toprule
    Конфигурация & Среднее \\
    \midrule
    Промптирование: prefix                       & 68,5 \\
    Промптирование: instruct                     & 69,3 \\
    Без этапа~1 (контрастивное предобучение)      & 68,7 \\
    Полный конвейер (три этапа)                   & 69,3 \\
    Полная модель ($\approx$3{,}0B)              & 69,3 \\
    Усечённая модель ($\approx$2{,}5B)           & 69,1 \\
    \bottomrule
  \end{tabular}
\end{table}

\noindent Абляции выполнены в условиях базовой публикации~[\ref{ref:acl}]
(средний балл по 23 задачам ruMTEB); абсолютный уровень отличается от
актуального публичного рейтинга (табл.~\ref{tab:res}) из-за последующего
обновления данных и версии бенчмарка, однако относительный вклад компонентов
сохраняется. Прунинг 25\,\% блоков снижает средний балл лишь на 0{,}2 при
сокращении числа параметров примерно на 17\,\%, что подтверждает низкую
значимость глубоких слоёв декодера для качества эмбеддингов и оправдывает
применение усечённой модели в развёртывании (раздел~\ref{sec:deploy}).

\subsection{Оценка на английской и кодовой дорожках MTEB}
Помимо русскоязычной дорожки, модель оценена на английской дорожке
MTEB(eng,~v2) и кодовой дорожке MTEB(code,~v1); сравнение с сильными базовыми
моделями приведено в табл.~\ref{tab:eng} и~\ref{tab:code}. При умеренном размере
(3,23~млрд параметров) сравнение проводится с моделями сопоставимого размера~---
до 4~млрд параметров; для предложенной модели приведена обновлённая версия
(раздел~\ref{sec:v2}). Баллы базовых моделей взяты из публичного рейтинга MTEB
для соответствующего набора задач; результаты обновлённой версии получены
авторами по тому же набору задач и протоколу агрегирования (среднее по задачам).

\begin{table}[ht]
  \centering
  \caption{Сравнение на MTEB(eng, v2) среди моделей до 4~млрд параметров
  (средний балл по задачам). Для предложенной модели приведена обновлённая
  версия (раздел~\ref{sec:v2})}\label{tab:eng}
  \begin{tabular}{lcc}
    \toprule
    Модель & Параметры, млрд & Mean (Task) \\
    \midrule
    \textbf{Giga-Embeddings (v2)}                    & \textbf{3,0}  & \textbf{71,93} \\
    jinaai/jina-embeddings-v5-text-small             & 0,60 & 71,78 \\
    IEITYuan/Yuan-embedding-2.0-en                   & 0,60 & 71,74 \\
    codefuse-ai/F2LLM-v2-1.7B                        & 1,72 & 71,63 \\
    NovaSearch/jasper\_en\_vision\_language\_v1      & 1,57 & 71,41 \\
    KaLM-embedding-multilingual-mini-instruct-v2.5   & 0,49 & 71,29 \\
    Giga-Embeddings-instruct (v1, лидерборд)         & 3,23 & 71,07 \\
    Qwen3-Embedding-0.6B                             & 0,60 & 70,47 \\
    google/embeddinggemma-300m                       & 0,31 & 69,67 \\
    multilingual-e5-large-instruct                   & 0,56 & 65,53 \\
    \bottomrule
  \end{tabular}
\end{table}

\begin{table}[ht]
  \centering
  \caption{Сравнение на MTEB(code, v1) среди моделей до 4~млрд параметров
  (средний балл). Для предложенной модели приведена обновлённая версия
  (раздел~\ref{sec:v2}, собственное измерение; базовая версия~--- 64,79)}\label{tab:code}
  \begin{tabular}{lcc}
    \toprule
    Модель & Параметры, млрд & Mean (Task) \\
    \midrule
    \textbf{Giga-Embeddings (v2)}    & \textbf{3,0}  & \textbf{79,30} \\
    codefuse-ai/F2LLM-v2-1.7B        & 1,72 & 78,76 \\
    codefuse-ai/F2LLM-v2-0.6B        & 0,60 & 77,41 \\
    perplexity-ai/pplx-embed-v1-0.6b & 0,60 & 75,85 \\
    codefuse-ai/F2LLM-v2-330M        & 0,33 & 75,74 \\
    codefuse-ai/C2LLM-0.5B           & 0,50 & 75,46 \\
    Qwen3-Embedding-0.6B             & 0,60 & 75,42 \\
    lightonai/LateOn-Code            & 0,15 & 74,12 \\
    Shuu12121/NightOwl-CodeEmbedding & 0,15 & 70,91 \\
    google/embeddinggemma-300m       & 0,31 & 68,76 \\
    Giga-Embeddings-instruct (v1, базовая) & 3,23 & 64,79 \\
    \bottomrule
  \end{tabular}
\end{table}

%% ===============================================================
\section{Анализ эффективности промышленного развёртывания}\label{sec:deploy}
%% ===============================================================
Усечённая модель ($\approx$2{,}5B) предназначена для снижения задержки и
ресурсопотребления при пренебрежимом снижении качества (потеря 0{,}2 балла,
табл.~\ref{tab:abl}). Оценим выигрыш аналитически. Прямой проход декодера
доминируется блоками трансформера (самовнимание и feed-forward), вычислительная
стоимость которых приблизительно равномерна по слоям. Удаление 9 из 36 блоков
сокращает число блоков на 25\,\%, что даёт примерно пропорциональное снижение
числа операций прямого прохода и объёма промежуточных активаций (включая
кэш ключей и значений), приходящихся на backbone. Поскольку таблица эмбеддингов
токенов и голова пулинга не затрагиваются, общее число параметров снижается на
меньшую величину~--- около 17\,\% (раздел~\ref{sec:arch}); соответственно
сокращается и объём памяти под веса.

Таким образом, усечение переводит примерно четверть вычислений backbone в
экономию задержки и памяти ценой 0{,}2 балла качества~--- благоприятный
компромисс для первичного этапа извлечения, работающего под жёсткими
ограничениями по задержке. Абсолютные значения задержки (мс/запрос), пропускной
способности (запросов/с) и резидентной памяти зависят от целевого оборудования,
длины входной последовательности, размера батча и движка инференса; их сквозное
измерение на конкретной конфигурации развёртывания составляет предмет инженерной
настройки.

%% ===============================================================
\section{Развитие модели: упрощение архитектуры и слияние весов}\label{sec:v2}
%% ===============================================================
Описанная модель получила дальнейшее развитие (далее~--- версия~v2),
направленное на упрощение архитектуры и повышение качества. Упрощение
архитектуры выполнено по запросам сообщества разработчиков открытого
программного обеспечения~--- для облегчения инференса и поддержки
специализированных движков; повышение качества достигнуто за счёт слияния весов
(model merging).

\subsection{Упрощение архитектуры}
Внесены два изменения. Во-первых, базовая архитектура заменена с архитектуры
семейства GigaChat на архитектуру, близкую к Qwen3, с двунаправленной маской
внимания. Во-вторых, удалён слой латентного пулинга внимания
(раздел~\ref{sec:arch}); это сокращает число параметров приблизительно на
10\,\%, упрощает граф вычислений и повышает совместимость модели с движками
быстрого инференса при сохранении качества.

\subsection{Обучение со слиянием весов}
Ключевым элементом обновлённого конвейера является слияние весов чекпоинтов,
обученных на различных датамиксах. Слияние выполняется методом сферической
линейной интерполяции (SLERP)~[\ref{ref:slerp}], сохраняющей норму
интерполируемых векторов параметров при усреднении моделей~[\ref{ref:mergekit}].
Конвейер состоит из шести этапов:
\begin{enumerate}
  \item три независимых предобучения (\textit{pretrain}) на различных датамиксах;
  \item слияние полученных предобученных чекпоинтов в единый;
  \item четыре независимых дообучения (\textit{fine-tuning}) слитого чекпоинта;
  \item слияние полученных дообученных чекпоинтов;
  \item мультизадачное дообучение (\textit{multitask});
  \item финальное слияние мультизадачного чекпоинта с дообученным.
\end{enumerate}
Схема конвейера приведена на рис.~\ref{fig:merge}. Усреднение весов моделей,
обученных на различных распределениях данных, повышает робастность представлений
и итоговое качество без увеличения стоимости инференса.

\begin{figure}[ht]
  \centering
  \resizebox{\textwidth}{!}{%
    \begin{tikzpicture}[
    font=\sffamily,
    node distance=0.7cm and 1.1cm,
    ck/.style={
        draw,
        rounded corners,
        minimum width=2.0cm,
        minimum height=0.9cm,
        align=center,
        thick
    },
    merge/.style={
        draw,
        diamond,
        aspect=2.2,
        inner sep=1pt,
        minimum width=2.0cm,
        align=center,
        thick
    },
    arr/.style={-{Stealth[length=3mm]}, thick}
]
    %% Этап 1: три предобучения
    \node[ck] (p1) {pretrain 1};
    \node[ck, below=of p1] (p2) {pretrain 2};
    \node[ck, below=of p2] (p3) {pretrain 3};
    %% Слияние 1
    \node[merge, right=of p2] (m1) {SLERP};
    %% Этап 3: четыре дообучения
    \node[ck, right=1.3cm of m1] (f1) {fine-tune 1};
    \node[ck, below=0.3cm of f1] (f2) {fine-tune 2};
    \node[ck, below=0.3cm of f2] (f3) {fine-tune 3};
    \node[ck, below=0.3cm of f3] (f4) {fine-tune 4};
    %% Слияние 2
    \node[merge, right=of f2] (m2) at ($(f2)!0.5!(f3) + (2.2cm,0)$) {SLERP};
    %% Мультизадачное + финальное слияние
    \node[ck, right=1.1cm of m2] (mt) {multitask};
    \node[merge, right=of mt] (m3) {SLERP};
    \node[ck, right=of m3] (out) {v2};

    %% Стрелки
    \draw[arr] (p1) -- (m1);
    \draw[arr] (p2) -- (m1);
    \draw[arr] (p3) -- (m1);
    \draw[arr] (m1) -- (f1);
    \draw[arr] (m1) -- (f2);
    \draw[arr] (m1) -- (f3);
    \draw[arr] (m1) -- (f4);
    \draw[arr] (f1) -- (m2);
    \draw[arr] (f2) -- (m2);
    \draw[arr] (f3) -- (m2);
    \draw[arr] (f4) -- (m2);
    \draw[arr] (m2) -- (mt);
    \draw[arr] (m2.south) to[out=-40,in=180] (m3.west);
    \draw[arr] (mt) -- (m3);
    \draw[arr] (m3) -- (out);
\end{tikzpicture}
}
  \caption{Конвейер обучения версии~v2 со слиянием весов: предобучения и
  дообучения на различных датамиксах объединяются методом SLERP на трёх уровнях
  (после предобучения, после дообучения и финально с мультизадачным
  чекпоинтом).}\label{fig:merge}
\end{figure}

\subsection{Результаты}
Обновлённая модель оценена на четырёх дорожках бенчмарка MTEB
(табл.~\ref{tab:v2}). Наибольший прирост достигнут на дорожке программного
кода~(+14{,}51). На русскоязычной дорожке модель сохраняет первое место. Среди
моделей до 4~млрд параметров (табл.~\ref{tab:eng} и~\ref{tab:code}) обновлённая
модель входит в число лучших на английской дорожке и занимает первое место на
дорожке программного кода. Результаты по многоязычной дорожке приводятся как
кросс-версионное сравнение (v1 против v2); сопоставление с базовыми моделями на
этой дорожке составляет предмет дальнейшей работы.

\begin{table}[ht]
  \centering
  \caption{Качество обновлённой модели на дорожках MTEB (средний балл)}\label{tab:v2}
  \begin{tabular}{lccc}
    \toprule
    Дорожка MTEB & Версия v1 (лидерборд) & Версия v2 & $\Delta$ \\
    \midrule
    Русский язык (rus)          & 74,16 & \textbf{74,57} & $+0{,}41$ \\
    Английский язык (eng)       & 71,07 & \textbf{71,93} & $+0{,}86$ \\
    Программный код (code)      & 64,79 & \textbf{79,30} & $+14{,}51$ \\
    Многоязычная (multilingual) & 55,51 & \textbf{60,06} & $+4{,}55$ \\
    \bottomrule
  \end{tabular}
\end{table}

\subsection{Масштабирование пайплайна}
Дальнейшее развитие пайплайна версии~v2 направлено на масштабирование трёх его
компонентов: (i)~более широкое слияние весов методом SLERP по большему числу
чекпоинтов, обученных на разнообразных датамиксах; (ii)~переход к архитектуре
Mixture of Experts с увеличенным числом экспертов; (iii)~расширение объёма и
разнообразия обучающих данных.

\subsection{Дальнейшая работа}
По описанному конвейеру в настоящее время обучается модель на 10~млрд параметров
с архитектурой Mixture of Experts; её оценка составляет предмет дальнейшей
работы. Планируется вклад в открытые движки инференса (vLLM, SGLang) для
упрощения развёртывания модели сообществом разработчиков открытого программного
обеспечения.

%% ===============================================================
\section{Заключение}
%% ===============================================================
Представлен фреймворк GigaEmbeddings~--- трёхэтапное инструктивное дообучение
декодерной языковой модели для русскоязычных текстовых эмбеддингов. Модель
достигает первого места на MTEB(rus, v1.1) (74,16), превосходя существенно более
крупные базовые модели; абляции подтверждают вклад контрастивного предобучения,
инструктивного промптирования и эффективность прунинга. Формализована интеграция
модели в двухэтапный поисковый конвейер с приближённым поиском ближайших соседей.
Дальнейшее развитие модели~--- упрощение архитектуры и обучение со слиянием весов
(раздел~\ref{sec:v2})~--- повышает качество (до 74,57 на русскоязычной дорожке) и
распространяет модель на английскую, кодовую и многоязычную дорожки MTEB; по тому
же конвейеру обучается модель на 10~млрд параметров с архитектурой Mixture of
Experts. Направление дальнейшей
работы~--- сквозная оценка качества и задержки полного поискового контура в
промышленных условиях.

%% ===============================================================
\section*{Список литературы}
%% ГОСТ Р 7.0.5-2008. Нумерация по порядку цитирования.
\begin{enumerate}
  \item\label{ref:acl} Kolodin~E., Khomich~D., Savushkin~N., Ianina~A., Minkin~F. GigaEmbeddings~--- Efficient Russian Language Embedding Model // Proceedings of the 10th Workshop on Slavic Natural Language Processing (Slavic NLP 2025). --- Association for Computational Linguistics, 2025. --- P.~17--24.
  \item\label{ref:gigachat} GigaChat team. GigaChat Family: Efficient Russian Language Modeling Through Mixture of Experts Architecture // Proceedings of the 63rd Annual Meeting of the Association for Computational Linguistics (Volume~3: System Demonstrations). --- Association for Computational Linguistics, 2025. --- P.~93--106.
  \item\label{ref:gromov} Gromov~A., Tirumala~K., Shapourian~H. et al. The Unreasonable Ineffectiveness of the Deeper Layers. --- 2024. --- arXiv:2403.17887.
  \item\label{ref:nvembed} Lee~C., Roy~R., Xu~M. et al. NV-Embed: Improved Techniques for Training LLMs as Generalist Embedding Models. --- 2024. --- arXiv:2405.17428.
  \item\label{ref:mteb} Snegirev~A., Tikhonova~M., Maksimova~A. et al. The Russian-Focused Embedders' Exploration: ruMTEB Benchmark and Russian Embedding Model Design. --- 2024. --- arXiv:2408.12503.
  \item\label{ref:qwen} Zhang~Y., Li~M., Long~D. et al. Qwen3 Embedding: Advancing Text Embedding and Reranking Through Foundation Models. --- 2025. --- arXiv:2506.05176.
  \item\label{ref:e5mistral} Wang~L., Yang~N., Huang~X. et al. Improving Text Embeddings with Large Language Models. --- 2024. --- arXiv:2401.00368.
  \item\label{ref:infonce} van den Oord~A., Li~Y., Vinyals~O. Representation Learning with Contrastive Predictive Coding. --- 2018. --- arXiv:1807.03748.
  \item\label{ref:s2orc} Lo~K., Wang~L.\,L., Neumann~M., Kinney~R., Weld~D. S2ORC: The Semantic Scholar Open Research Corpus // Proceedings of the 58th Annual Meeting of the Association for Computational Linguistics. --- 2020. --- P.~4969--4983.
  \item\label{ref:msmarco} Bajaj~P., Campos~D., Craswell~N. et al. MS MARCO: A Human Generated MAchine Reading COmprehension Dataset. --- 2018. --- arXiv:1611.09268.
  \item\label{ref:nq} Kwiatkowski~T., Palomaki~J., Redfield~O. et al. Natural Questions: A Benchmark for Question Answering Research // Transactions of the Association for Computational Linguistics. --- 2019. --- Vol.~7. --- P.~452--466.
  \item\label{ref:squad} Rajpurkar~P., Zhang~J., Lopyrev~K., Liang~P. SQuAD: 100,000+ Questions for Machine Comprehension of Text // Proceedings of the 2016 Conference on Empirical Methods in Natural Language Processing. --- 2016. --- P.~2383--2392.
  \item\label{ref:miracl} Zhang~X., Thakur~N., Ogundepo~O. et al. MIRACL: A Multilingual Retrieval Dataset Covering 18 Diverse Languages // Transactions of the Association for Computational Linguistics. --- 2023. --- Vol.~11. --- P.~1114--1131.
  \item\label{ref:mrtydi} Zhang~X., Ma~X., Shi~P., Lin~J. Mr.~TyDi: A Multi-lingual Benchmark for Dense Retrieval // Proceedings of the 1st Workshop on Multilingual Representation Learning. --- 2021. --- P.~127--137.
  \item\label{ref:hnsw} Malkov~Yu.\,A., Yashunin~D.\,A. Efficient and Robust Approximate Nearest Neighbor Search Using Hierarchical Navigable Small World Graphs // IEEE Transactions on Pattern Analysis and Machine Intelligence. --- 2020. --- Vol.~42, no.~4. --- P.~824--836.
  \item\label{ref:faiss} Johnson~J., Douze~M., J\'egou~H. Billion-Scale Similarity Search with GPUs // IEEE Transactions on Big Data. --- 2021. --- Vol.~7, no.~3. --- P.~535--547.
  \item\label{ref:bgem3} Chen~J., Xiao~S., Zhang~P. et al. BGE M3-Embedding: Multi-Lingual, Multi-Functionality, Multi-Granularity Text Embeddings Through Self-Knowledge Distillation. --- 2024. --- arXiv:2402.03216.
  \item\label{ref:gritlm} Muennighoff~N., Su~H., Wang~L. et al. Generative Representational Instruction Tuning. --- 2024. --- arXiv:2402.09906.
  \item\label{ref:slerp} Shoemake~K. Animating Rotation with Quaternion Curves // ACM SIGGRAPH Computer Graphics. --- 1985. --- Vol.~19, no.~3. --- P.~245--254.
  \item\label{ref:mergekit} Goddard~C., Siriwardhana~S., Ehghaghi~M. et al. Arcee's MergeKit: A Toolkit for Merging Large Language Models. --- 2024. --- arXiv:2403.13257.
\end{enumerate}

\medskip
\noindent\textbf{Сведения об авторах.}
Колодин Егор Ильич~--- аспирант, Московский физико-технический институт
(национальный исследовательский университет). \\
Янина Анастасия Олеговна~--- кандидат физико-математических наук, научный
сотрудник лаборатории фундаментальных исследований искусственного интеллекта,
Московский физико-технический институт (национальный исследовательский
университет).

%% ---------------------------------------------------------------
%%  АНГЛИЙСКАЯ ЧАСТЬ (структура дублирует русскую точным переводом)
%%  [ЗАПОЛНИТЬ: полный перевод тела статьи. Ниже — front matter.]
%% ---------------------------------------------------------------
\selectlanguage{english}
\bigskip
\begin{center}
  {\large\bfseries Hierarchical Instruction Tuning of a Decoder Language Model\\
  for Russian Text Embeddings}

  \medskip
  {\bfseries E.\,I.~Kolodin$^{1}$, A.\,O.~Yanina$^{1}$}

  \smallskip
  {\itshape $^{1}$ Moscow Institute of Physics and Technology (National Research University)}
\end{center}

\medskip
\noindent\textbf{Abstract.} % exact translation of the Russian abstract, 100-120 words
We present GigaEmbeddings, a framework for high-performance Russian text
embeddings built on a decoder language model. A hierarchical three-stage
instruction-tuning pipeline is proposed: adapting the decoder to bidirectional
attention and latent attention pooling, contrastive pre-training on web-scale
corpora, fine-tuning with hard negatives, and multi-task generalization. We
formalize the model's integration into a two-stage retrieval pipeline with
approximate nearest-neighbor search and analyze the efficiency of deep-layer
pruning for production deployment. On the MTEB(rus, v1.1) benchmark the model
ranks first with a mean score of 74.16, outperforming substantially larger
multilingual and Russian baselines. Ablations confirm the contribution of each
training stage.

\medskip
\noindent\textbf{Keywords:} text embeddings, information retrieval, large language
models, contrastive learning, instruction tuning, approximate nearest-neighbor
search, Russian language, MTEB.

%% [ЗАПОЛНИТЬ: полный английский перевод разделов 1–7 и References,
%%  идентичных по нумерации русскому списку литературы.]
\selectlanguage{russian}

\end{document}
